\documentclass[UTF8,a4paper,12pt,fontset=fandol]{ctexart}
\usepackage[margin=2.5cm]{geometry}
\usepackage{array,tabularx,booktabs,enumitem,hyperref,fancyhdr}
\hypersetup{hidelinks}
\pagestyle{fancy}\fancyhf{}\fancyfoot[C]{\thepage}\renewcommand{\headrulewidth}{0pt}
\linespread{1.15}\setlength{\parindent}{2em}\setlist{nosep}
\newcommand{\blank}{\underline{\hspace{3cm}}}
\newcolumntype{Y}{>{\raggedright\arraybackslash}X}
\begin{document}
\begin{center}{\zihao{2}\bfseries AI工具使用详情}\par\vspace{8pt}论文支撑材料\end{center}
\noindent 本文件记录AI辅助范围。工具与使用时间、采纳内容和人工处理情况应由参赛队据实补充，空栏表示尚未填写。

\section*{一、使用工具}
\par\noindent\begin{tabularx}{\linewidth}{p{.28\linewidth}Y}\toprule
项目 & 记录\\\midrule
工具 & OpenAI Codex\\
模型信息 & 当前任务系统标识为基于GPT-6的编程代理；其他使用情况待补：\blank\\
使用日期 & 2026年9月13日；完整使用时段待补：\blank\\
其他AI工具 & 工具、模型及日期：\blank\\\bottomrule
\end{tabularx}

\section*{二、辅助范围}
\begin{enumerate}
\item \textbf{建模与程序。} 讨论搜索布局、测向几何、定向接收、任务排序及光学清除，辅助编写、调试和整理Python程序。
\item \textbf{数据与图表。} 辅助运行合成案例，整理已有计算结果和演练截图数据，绘制几何示意与费用图。合成案例不作为官方测试成绩。
\item \textbf{论文与附录。} 辅助撰写摘要、模型推导、算法描述、结果分析和局限，编写LaTeX排版、完整程序附录及打包程序。未取得的实测数据保留空栏。
\item \textbf{材料处理。} 依据所提供的题目、接口说明、写作流程和排版资料组织材料。图表与正文引用来源需由参赛队实际阅读确认，结构借鉴不包括复制他人正文或数据。
\end{enumerate}

\section*{三、代表性实际请求}
\par\noindent\begin{tabularx}{\linewidth}{p{.49\linewidth}Y}\toprule
请求摘录 & 辅助处理\\\midrule
“要求方法能确保找到所有干扰点” & 采用完整搜索与定位流程，说明适用条件和终止条件。\\
“可以 进行代码的书写” & 辅助实现求解程序与运行入口。\\
“请你提出几个最有可能的可行的方案，分别列出优缺点” & 讨论方法选择、成本与适用条件。\\
“帮我撰写一份30页的论文……还没有测的数据……先空出来” & 组织30页主稿及完整程序附录，未测项据实留空。\\\bottomrule
\end{tabularx}

\clearpage
\section*{四、实际采纳与人工处理（待填）}
下表由参赛队按实际情况补充，可注明采纳、修改或未采纳的内容。使用匿名代号及日期，不在评阅材料中填写姓名或学校。

\par\noindent\begin{tabularx}{\linewidth}{p{.35\linewidth}Y}\toprule
事项 & 实际处理、代号及日期\\\midrule
模型、推导和适用条件 & \blank\\[16pt]
程序理解与运行 & \blank\\[16pt]
计算结果及图表 & \blank\\[16pt]
演练和正式成绩 & \blank\\[16pt]
文献阅读与引用 & \blank\\[16pt]
论文正文及完整程序附录 & \blank\\[16pt]
其他AI辅助内容 & \blank\\[16pt]\bottomrule
\end{tabularx}

\noindent 工具生成的内容、自动运行结果与参赛队实际理解和确认的内容应分别记录。正式数据依据真实测试填写；不得将合成结果或空栏写成已完成的正式成绩。
\end{document}
